Bug reports provide the information developers need to triage and resolve software defects, yet unclear or unstructured reports delay this process. Existing work on Qwen2.5-7B consistently relies on supervised fine-tuning or reinforcement learning before evaluation, while GPT-4o serves only as a few-shot baseline; no study has examined whether Qwen2.5-7B under few-shot prompting only -- without additional training -- can match the GPT-4o thresholds established by Acharya and Ginde (2025) (CTQRS $\geq$75\%, ROUGE-1 $\geq$0.47, SBERT $\geq$0.82). We evaluate Qwen2.5-7B-Instruct with 3-shot prompting on 262 Bugzilla bug reports, measuring CTQRS, ROUGE-1 F1, and SBERT cosine similarity against these thresholds using one-sample Wilcoxon signed-rank tests. Qwen2.5-7B exceeds the CTQRS threshold (median=88.89\%, $p<0.001$, Cliff's $\delta$=0.908) and the ROUGE-1 threshold (median=0.491, $p$=0.002, $\delta$=0.191), but falls short of the SBERT threshold (median=0.632, $p$=1.000, $\delta$=-0.969); a length-constrained ablation raises ROUGE-1 to 0.532 but SBERT remains at 0.609, indicating stylistic verbosity rather than semantic inaccuracy. These results show that few-shot prompting alone lets an open-source LLM match GPT-4o's structural and lexical quality without additional training, but output style remains the key barrier to semantic-similarity metrics.